
The primal problem is
\[
\min_{x \in \mathbb{R}^2} 4x_1+2x_2
\]
subject to 
\begin{align*}
x_1+x_2 & \geq 2,\\
-x_1-x_2 & \geq 1,\\
x_1 & \geq 0,\\
x_2 & \geq 0.
\end{align*}
The quantity $x_1+x_2$ must be both larger or equal to 2, and lesser
or equal to -1. This is not possible, and the problem is
infeasible. Consequently, we expect the dual problem to be either
infeasible or unbounded.

  To write the Lagrangian, we first write the problem as follows:
$$\min_{x \in \mathbb{R}^2} 4x_1+2x_2$$ 
subject to 
\begin{alignat*}{5}
& g_1(x) && = 2 && -x_1 && -x_2 && \leq 0 \quad (\mu_1),\\
& g_2(x) &&= 1 && +x_1 && +x_2 && \leq 0 \quad (\mu_2),\\
& g_3(x) &&= && -x_1 && && \leq 0 \quad (\mu_3),\\
& g_4(x) &&= && && -x_2 && \leq 0 \quad (\mu_4),
\end{alignat*}
where $\mu_1$, $\mu_2$, $\mu_3$ and $\mu_3$ are the penalty parameters
associated with each constraint.

The Lagrangian function is 
\begin{align*}
L(x_1,x_2,\mu_1,\mu_2,\mu_3,\mu_4) & = 4x_1+2x_2 + \mu_1(2-x_1-x_2) \\
& \qquad + \mu_2(1 +x_1+x_2) - \mu_3 x_1 - \mu_4 x_2\\
& = (4-\mu_1+\mu_2-\mu_3)x_1\\ & \qquad + (2-\mu_1+\mu_2-\mu_4)x_2 + 2\mu_1 + \mu_2.
\end{align*}

The dual function is
\[
q(\mu_1,\mu_2,\mu_3,\mu_4) = \min_{x_1,x_2} L(x_1,x_2,\mu_1,\mu_2,\mu_3,\mu_4).
\]

In order for the dual function to be bounded, the coefficients of
$x_1$ and $x_2$ have to be zero.  Therefore, we have to impose the
following constraints on the penalty parameters:
\begin{align}
  \mu_3&=4-\mu_1+\mu_2, \label{eq:5} \\
  \mu_4&=2-\mu_1+\mu_2. \label{eq:6}
\end{align}
In that case, the dual function becomes 
\[
q(\mu_1,\mu_2,\mu_3,\mu_4) = 2\mu_1 + \mu_2.
\]
Moreover, the parameters associated with the inequality constraints
must be non negative. That is, $\mu_1 \geq 0$, $\mu_2 \geq
0$, $\mu_3 \geq 0$, $\mu_4 \geq 0$. The dual problem is therefore
\[
\max_{\mu_1,\mu_2,\mu_3,\mu_4}2\mu_1 + \mu_2
\]
subject to
\begin{align*}
  \mu_3&=4-\mu_1+\mu_2,\\
  \mu_4&=2-\mu_1+\mu_2, \\
  \mu_1 & \geq 0, \\
  \mu_2 & \geq 0, \\
  \mu_3 & \geq 0, \\
  \mu_4 & \geq 0.
\end{align*}

As the parameters $\mu_3$ and $\mu_4$ do not appear in the objective
function, we 
can rewrite
\req{eq:5} and \req{eq:6} as
\begin{align*}
  4-\mu_1+\mu_2 & \geq 0,\\
  2-\mu_1+\mu_2 & \geq 0, \\
\end{align*}
and eliminate $\mu_3$ and $\mu_4$, so that the dual problem becomes
\[
\max_{\mu_1,\mu_2}2\mu_1 + \mu_2
\]
subject to
\begin{align*}
  4-\mu_1+\mu_2 & \geq 0,\\
  2-\mu_1+\mu_2 & \geq 0, \\
  \mu_1 & \geq 0, \\
  \mu_2 & \geq 0.
\end{align*}
The problem happens to be unbounded. To see it, first fix the value of $\mu_1$ to
0. We obtain
\[
\max_{\mu_2} \mu_2
\]
subject to
\begin{align*}
  \mu_2 & \geq -4,\\
  \mu_2 & \geq -2, \\
  \mu_2 & \geq 0.
\end{align*}
As $\mu_2$ is not constrained from above, its value and, therefore,
the value of the objective function, can always be made larger than any
arbitrary value. 

